\chapter{三相记忆的软件几何实现：GPU上的What--Where双通路与Soft-HGU}

\section{从物理实现回到软件实现}

上一章讨论了三相记忆在连续物理动力介质中的可能实现：
\[
h
\leftrightarrow
E
\leftrightarrow
\Theta,
\]
其中 $h$ 表示当前仍然活跃的认知状态，$E$ 表示已经发生但仍保持事件连续性的经历结构，$\Theta$ 表示通过长期学习形成、能够约束和生成未来认知过程的慢结构。

这一结构本身并不要求一定存在专用物理芯片。它所要求的只是系统内部至少存在三组清晰分离的时间尺度：
\[
\tau_h\ll\tau_E\ll\tau_\Theta.
\]

因此，在专门物理介质尚不存在的条件下，可以利用现有 GPU 对同一动力结构进行软件模拟。此时 GPU 仍然执行矩阵乘法、张量运算与离散数值积分，但被模拟的对象不再必须是传统 Transformer 的逐层前向传播，而可以是一个具有内部状态、几何结构、时间尺度、残留状态和情景绑定机制的连续认知动力系统。

因此需要区分两个层面：
\[
\boxed{
\mathrm{ExecutionSubstrate}
\neq
\mathrm{CognitiveDynamics}.
}
\]

GPU 是执行介质，而软件中定义的状态空间、动力学、记忆相态和关系绑定机制才是认知系统本身。

由此定义
\[
\boxed{
\mathrm{SoftGeometricMemorySystem}
}
\]
为：利用现有代数计算硬件，对具有显式表示空间、状态动力学、多时间尺度记忆和关系绑定机制的几何认知系统进行数值模拟的软件架构。

它不是对物理芯片的简单接口模拟，而首先是一种可以独立运行、训练和验证的三相记忆实现。

\section{软件几何系统的最小状态定义}

为了构造软件三相记忆，首先定义一个最小运行状态
\[
\mathcal{X}(t)
=
\left(
X_{\mathrm{str}}(t),
X_{\mathrm{sem}}(t),
\Phi(t),
Z(t)
\right).
\]

其中：

\[
X_{\mathrm{str}}(t)
\]
表示当前结构、位置、拓扑或关系骨架状态；

\[
X_{\mathrm{sem}}(t)
\]
表示当前对象、内容、属性和语义状态；

\[
\Phi(t)
\]
表示用于描述不同局部状态之间时间关系、相位关系或动态同步状态的附加动力变量；

\[
Z(t)
\]
表示由当前大量局部状态压缩得到的较低维全局调制状态。

这里不要求这些变量必须具有某种特定神经科学含义。它们只是构造软件动力系统所需要的四类基本变量。

在具体实现中，可以把
\[
X_{\mathrm{str}}
\]
映射为当前代码中的结构流状态，将
\[
X_{\mathrm{sem}}
\]
映射为语义流状态，将 $\Phi$ 映射为相位和频率状态，而 $Z$ 则对应对整个当前认知场进行压缩后的宏观控制变量。

从三相记忆角度看，这些变量当前全部属于
\[
h(t),
\]
即
\[
\boxed{
h(t)
=
\mathcal{X}(t).
}
\]

因此，软件实现中的工作记忆首先不是一个独立的数组，而是系统在当前计算时刻仍然参与动力演化的完整状态集合。

\section{为什么软件几何系统需要显式双通路}

如果所有对象语义、空间位置、动作关系和时序结构都被提前压缩到单一向量中，那么系统虽然能够学习联合表示，却会逐渐失去一个重要能力：将相对稳定的对象身份与不断变化的现实位置进行自由重新组合。

例如，设对象为
\[
C_{\mathrm{cup}},
\]
而位置分别为
\[
S_1=\mathrm{kitchen},
\qquad
S_2=\mathrm{bedroom}.
\]

如果系统为每一种组合都学习一个独立状态
\[
C_{\mathrm{cup@kitchen}},
\qquad
C_{\mathrm{cup@bedroom}},
\]
那么对象数量和场景数量会共同扩大表示空间。

更加一般的做法是维持两个状态空间：
\[
\boxed{
\mathcal{W}
}
\]
和
\[
\boxed{
\mathcal{S}.
}
\]

其中 $\mathcal{W}$ 称为 What Space，主要承载对象身份、概念、属性、功能和语义关系；

$\mathcal{S}$ 称为 Where Space，主要承载空间位置、方向、距离、时间位置、运动结构以及与当前身体的相对关系。

于是当前状态分别表示为
\[
X_W(t)\in\mathcal{W},
\]
\[
X_S(t)\in\mathcal{S}.
\]

二者满足
\[
\boxed{
\mathrm{Identity}
\neq
\mathrm{Situation}.
}
\]

对象身份可以长期稳定，而对象所处场景持续变化。

因此软件双通路的基本目标不是复制生物脑结构，而是提供
\[
\boxed{
\mathrm{StableSemanticIdentity}
+
\mathrm{VariableSpatiotemporalContext}.
}
\]

\section{从结构流与语义流到What--Where表示}

若已有模型已经分别维护结构状态和语义状态，则不需要立即重构整个模型。更加稳妥的方法是在现有状态上增加两个显式投影：
\[
K_S
=
P_S(X_{\mathrm{str}}),
\]
\[
V_W
=
P_W(X_{\mathrm{sem}}),
\]
其中 $P_S$ 和 $P_W$ 是可学习投影。

这里首先使用符号
\[
K_S
\]
表示当前可用于情景寻址的 Where 表示，使用
\[
V_W
\]
表示当前可被情景记忆保存和恢复的 What 表示。

需要特别强调：
\[
\boxed{
X_{\mathrm{str}}
\neq
X_S
}
\]
以及
\[
\boxed{
X_{\mathrm{sem}}
\neq
X_W
}
\]
并非严格恒等。

原始结构流可以同时包含文本顺序、一般拓扑和物理空间等多种结构信息，因此真正进入情景记忆的 Where 状态最好通过专门投影产生：
\[
X_S
=
P_{\mathrm{where}}(X_{\mathrm{str}}).
\]

同理，原始语义流中可能包含大量短期计算信息，真正进入情景系统的 What 状态应由
\[
X_W
=
P_{\mathrm{what}}(X_{\mathrm{sem}})
\]
提取。

这样可以保持底层动力系统的一般性，同时逐渐形成明确的 What--Where 认知接口。

\section{What--Where双通路的动力学形式}

双通路不应该只是两个静态 embedding，而应该分别具有自己的状态演化：
\[
\frac{dX_W}{dt}
=
F_W
\left(
X_W,
C_{SW}(X_S),
u_W,
Z
\right),
\]
\[
\frac{dX_S}{dt}
=
F_S
\left(
X_S,
C_{WS}(X_W),
u_S,
Z
\right).
\]

其中 $F_W$ 和 $F_S$ 分别表示语义流与空间流的局部动力学；

$C_{SW}$ 表示空间状态对语义状态的有限调制；

$C_{WS}$ 表示语义状态对空间状态的有限调制；

$u_W,u_S$ 表示外部输入；

$Z$ 表示较慢的全局控制变量。

核心原则是
\[
\boxed{
\mathrm{Coupled}
\neq
\mathrm{Collapsed}.
}
\]

What 与 Where 可以相互影响，但不必在每一层中彻底融合为同一表示。

这样，同一个对象可以在不同地点被重新绑定，而同一个地点也可以承载不同对象和事件。

\section{双通路仍然不足以形成情景记忆}

即使系统已经能够同时表示
\[
X_W(t)
\]
和
\[
X_S(t),
\]
仍然只能回答
\[
\text{现在看到了什么}
\]
和
\[
\text{它现在在哪里}.
\]

真正的情景记忆还必须回答：
\[
\boxed{
\text{过去的什么对象，在什么时间、什么地方，发生了什么变化？}
}
\]

因此需要一个新的功能层：
\[
\boxed{
\mathrm{EpisodicBindingLayer}.
}
\]

这一层负责把当前分散存在于多个表示空间中的状态绑定成一个具有时间连续性的事件结构。

为了与未来可能的物理情景绑定单元区分，本章将这一软件组件定义为
\[
\boxed{
\mathrm{Soft\text{-}HGU}
}
\]
即 Software Hippocampal Geometric Unit。

这里的名称只表示它承担快速情景绑定、索引、模式补全和情景召回功能；其软件实现可以完全基于 GPU Tensor，不要求任何专用硬件。

\section{Soft-HGU的最小形式：Where-Key与What-Value}

Soft-HGU 的第一代实现可以非常简单。

定义每一个 Episode：
\[
E_i
=
(K_i,V_i,M_i),
\]
其中
\[
K_i
=
P_K(X_{S,i}),
\]
\[
V_i
=
P_V(X_{W,i}),
\]
而 $M_i$ 表示与该情景相关的元数据。

这里使用 Key--Value 语言并不意味着 Soft-HGU 等价于 Transformer 的 KV Cache。

其功能定义完全不同：

\[
\boxed{
K_i:
\text{在什么情境下寻找这段经历},
}
\]
\[
\boxed{
V_i:
\text{这段经历包含什么内容}.
}
\]

最简单地，
\[
K_i=\mathrm{Where}_i,
\]
\[
V_i=\mathrm{What}_i.
\]

于是，当 Agent 再次到达相似地点时，可以通过当前位置形成查询
\[
q_S=P_Q(X_S(t)),
\]
并计算
\[
r_i
=
\operatorname{sim}(q_S,K_i).
\]

选取相关 Episode 后：
\[
\hat V
=
\sum_i
\alpha_i V_i,
\]
其中
\[
\alpha_i
=
\frac{\exp(r_i)}
{\sum_j\exp(r_j)}.
\]

这样便形成：
\[
\boxed{
\mathrm{CurrentWhere}
\rightarrow
\mathrm{PastWhat}.
}
\]

例如：
\[
\mathrm{kitchen/table/right}
\rightarrow
\mathrm{blue\ cup}.
\]

这构成最简单的情景模式补全。

\section{为什么情景Key必须逐渐从Where扩展到Where--When}

仅仅使用空间作为 Key 会迅速产生冲突。

同一个地点在不同时间可能发生大量完全不同的事件。因此更合理的情景键为
\[
\boxed{
K_i
=
P_K
(
X_{S,i},
T_i
),
}
\]
其中 $T_i$ 表示事件发生时的时间结构。

由此
\[
\mathrm{kitchen@morning}
\]
和
\[
\mathrm{kitchen@night}
\]
可以成为不同情景索引。

进一步还可以加入主体上下文：
\[
\boxed{
K_i
=
P_K
(
X_{S,i},
T_i,
C^{\mathrm{self}}_i
).
}
\]

于是 Key 所表达的已经不是简单坐标，而是
\[
\boxed{
\mathrm{SpatiotemporalSelfContext}.
}
\]

这更加接近真正的情景记忆寻址问题。

\section{What也必须逐渐扩展为事件内容}

第一代系统可以令
\[
V_i=X_{W,i},
\]
即保存当时的对象和语义状态。

但是成熟的 Episode 不应该只包含 What，还需要包含：
\[
\mathrm{Action},
\qquad
\mathrm{Relation},
\qquad
\mathrm{Outcome},
\qquad
\mathrm{Affect}.
\]

因此更一般的 Value 应写成
\[
\boxed{
V_i
=
P_V
(
X_{W,i},
A_i,
R_i,
O_i,
F_i
),
}
\]
其中：

$A_i$ 表示该事件中的动作；

$R_i$ 表示对象之间的关系；

$O_i$ 表示结果；

$F_i$ 表示事件对主体产生的功能影响摘要。

于是完整 Episode 可以表示为
\[
\boxed{
E_i
=
(
K_i^{\mathrm{where/when/self}},
V_i^{\mathrm{what/action/relation/outcome}},
M_i
).
}
\]

此时 Key--Value 只是软件实现接口，真正的本体仍然是一个关系事件。

\section{情景记忆必须支持双向查询}

如果情景系统只能实现
\[
\mathrm{Where}\rightarrow\mathrm{What},
\]
那么用户问
\[
\text{蓝色杯子在哪里？}
\]
时就无法自然完成回忆。

因此成熟 Soft-HGU 应支持双向或多入口模式补全。

一种简单方案是同时建立：
\[
K_i^S=P_S(X_{S,i}),
\]
\[
K_i^W=P_W(X_{W,i}),
\]
并令 Episode 保存
\[
E_i
=
(
K_i^S,
K_i^W,
V_i,
M_i
).
\]

于是可以执行：
\[
\mathrm{Where}
\rightarrow
E_i
\rightarrow
\mathrm{What},
\]
也可以执行
\[
\mathrm{What}
\rightarrow
E_i
\rightarrow
\mathrm{Where}.
\]

从更加一般的角度看，Soft-HGU 不应被理解成单向字典：
\[
K\rightarrow V,
\]
而应理解成
\[
\boxed{
\mathrm{RelationalAssociativeMemory}.
}
\]

也就是说，只要给定 Episode 中足够一部分结构，系统就可以恢复剩余关系。

\section{工作记忆h在软件双通路系统中的定义}

在 Soft-HGU 加入以后，工作记忆需要重新写成分布式形式：
\[
\boxed{
h(t)
=
[
X_W(t),
X_S(t),
\Phi(t),
X_H^{\mathrm{active}}(t),
Z(t)
].
}
\]

其中 $X_H^{\mathrm{active}}$ 表示当前被召回或正在形成的 Episode 状态。

因此软件系统中的工作记忆不是中央 KV Cache，而是
\[
\boxed{
\mathrm{TemporarilyCoherentDistributedState}.
}
\]

也就是说，只有当前需要协同参与同一认知过程的 What、Where、Episode 和全局调制状态形成当前 $h$。

这与此前定义
\[
\boxed{
h=\mathrm{SparseExplicitCrossScaleSurface}
}
\]
完全兼容。

大量局部视觉状态、运动控制状态和低层预测过程可以继续留在局部系统中，并不需要全部复制进入 $h$。

因此：
\[
\boxed{
\mathrm{MostDynamics}\notin h.
}
\]

成熟智能的关键不是扩大 $h$ 到无限，而是只让真正需要跨系统协调的结构进入当前工作面。

\section{情景记忆E必须表示轨迹，而不只是状态快照}

三相记忆中的
\[
E=\mathrm{Trajectory}
\]
意味着情景记忆不能退化成一系列无关联的瞬时快照。

因此一个真正完整的软件 Episode 应同时描述事件开始、变化和结束。

可以定义：
\[
E_i
=
(
K_i,
V_i,
\Delta_i,
M_i
),
\]
其中
\[
\Delta_i
=
(
Z_i^{start},
Z_i^{peak},
Z_i^{end},
\Delta t_i,
O_i
).
\]

这里：

$Z_i^{start}$ 表示事件起始宏观状态；

$Z_i^{peak}$ 表示事件过程中最显著状态；

$Z_i^{end}$ 表示事件结束状态；

$\Delta t_i$ 表示持续时间；

$O_i$ 表示事件结果。

于是情景系统保存的不仅是
\[
\text{当时有什么},
\]
还保存
\[
\boxed{
\text{发生了什么变化}.
}
\]

因此
\[
E
\]
真正成为
\[
\boxed{
\mathrm{StateTransitionMemory}
}
\]
而不仅仅是对象检索数据库。

\section{事件边界：什么时候开始形成一个Episode}

情景记忆不能每一个计算 step 都写一条记录，否则情景数量将无限增长，而且真正事件结构会被大量连续微小变化淹没。

因此必须首先定义
\[
\boxed{
\mathrm{EpisodeBoundary}.
}
\]

设当前运行状态变化量为
\[
\delta_t
=
\|h(t)-h(t-\Delta t)\|.
\]

可以观察到，一个典型事件往往具有如下动力过程：
\[
\boxed{
\mathrm{StableState}_A
\rightarrow
\mathrm{Transition}
\rightarrow
\mathrm{StableState}_B.
}
\]

因此可以定义：

当
\[
\delta_t>\theta_{\mathrm{start}}
\]
且状态持续偏离原吸引子时，开始记录事件；

当
\[
\delta_t<\theta_{\mathrm{end}}
\]
并保持一段时间以后，结束 Episode。

这样得到：
\[
t_i^{start},
\qquad
t_i^{end}.
\]

完整 Episode 则由这一区间上的轨迹摘要形成：
\[
E_i
=
\mathcal{B}
\left(
h[t_i^{start}:t_i^{end}]
\right).
\]

这里 $\mathcal{B}$ 表示情景绑定与压缩算子。

因此 Episode 不是按照固定 token 长度切割，而是按照系统动力学自身的变化边界形成。

\section{从动态弛豫信号构造事件边界}

如果软件动力系统已经具有弛豫检测机制，则无需重新设计事件边界估计器。

设局部认知应力为
\[
\sigma_c(t)
=
\|
X_{\mathrm{sem}}^{new}
-
X_{\mathrm{sem}}^{old}
\|.
\]

当
\[
\sigma_c(t)
\]
长期较低时，可以认为局部状态已经进入相对稳定区域。

而当
\[
\sigma_c(t)
\]
突然升高，则意味着当前结构正在经历显著变化。

因此可以定义：
\[
\mathrm{EpisodeStart}
:
\sigma_c
\uparrow
\text{ and }
\sigma_c>\theta_h,
\]
\[
\mathrm{EpisodeEnd}
:
\sigma_c
\downarrow
\text{ and }
\sigma_c<\theta_l.
\]

其中
\[
\theta_h>\theta_l
\]
形成迟滞区间，以避免状态在阈值附近反复开启和关闭 Episode。

于是事件切分本身也成为一个动力学检测问题，而不是外部人工标注问题。

\section{为什么不是所有Episode都应该进入E}

即使已经识别出事件边界，也不意味着所有 Episode 都值得长期保存。

三相记忆必须解决
\[
h\rightarrow E
\]
的选择问题。

定义情景写入得分
\[
s_i^{write}
=
\alpha N_i
+
\beta R_i
+
\gamma G_i
+
\delta S_i
+
\eta A_i
-
\lambda C_i,
\]
其中：

$N_i$ 表示新颖度；

$R_i$ 表示预测残差；

$G_i$ 表示目标相关度；

$S_i$ 表示自我相关度；

$A_i$ 表示内稳态或价值显著性；

$C_i$ 表示记忆写入和未来维护成本。

写入概率定义为
\[
p_i^{write}
=
\sigma
\left(
s_i^{write}
\right).
\]

只有满足
\[
p_i^{write}>\theta_E
\]
时，Episode 才真正进入情景记忆集合。

因此：
\[
\boxed{
\mathrm{EventDetected}
\neq
\mathrm{EpisodeStored}.
}
\]

一个系统可以经历大量现实变化，却只选择少数具有长期认知价值的经历成为自己的情景历史。

\section{全局调制状态如何参与情景写入}

如果软件系统已经存在一个从当前分布式状态压缩形成的全局变量
\[
Z(t),
\]
则可以直接把它作为第一代情景写入控制的输入。

定义：
\[
Z(t)
=
\mathcal{C}
\left(
X_W,
X_S,
\Phi,
X_H^{active}
\right),
\]
其中 $\mathcal{C}$ 为全局压缩算子。

然后：
\[
p_i^{write}
=
W_E
\left(
Z_i,
\sigma_{c,i},
H_i,
N_i
\right),
\]
其中 $H_i$ 可以表示当前内部路由、选择或竞争的熵。

这一设计使系统不需要在第一版立即实现完整自我系统。它首先只需要学会：
\[
\boxed{
\text{当前发生的变化是否值得留下？}
}
\]

以后再逐步把 Self、Goal、Affect 等更慢结构加入同一写入函数。

\section{情景记忆的流变：E不是静态数据库}

三相记忆中的 $E$ 被定义为一种比 $h$ 稳定、比 $\Theta$ 更可塑的中时间尺度结构。

因此情景记忆本身必须能够衰减、强化、合并和消失。

为每一个 Episode 定义记忆强度
\[
m_i(t).
\]

其演化可以写为
\[
\boxed{
\tau_E\dot m_i
=
-\lambda_i m_i
+
\alpha_r R_i(t)
+
\alpha_v V_i
+
\alpha_u U_i(t),
}
\]
其中：

$\lambda_i$ 表示自然衰减；

$R_i(t)$ 表示该 Episode 是否被重新召回；

$V_i$ 表示长期价值；

$U_i(t)$ 表示后续经历是否不断重新支持该事件。

如果长期满足
\[
m_i(t)\rightarrow0,
\]
则 Episode 可以被压缩、合并或删除。

如果
\[
m_i(t)\uparrow,
\]
则它可以成为长期巩固候选。

因此：
\[
\boxed{
E
=
\mathrm{DynamicEpisodePopulation}.
}
\]

情景记忆不是不断增长的日志，而是一个持续演化的事件群体。

\section{情景召回不是普通向量搜索}

最简单实现可以直接使用相似度：
\[
r_i=q^\top K_i.
\]

但是在双通路架构中，可以进一步把召回拆成多个匹配因素：
\[
\boxed{
R_i
=
R_S(i)
R_W(i)
R_T(i)
R_G(i)
R_\Psi(i).
}
\]

其中
\[
R_S(i)
\]
表示空间匹配；

\[
R_W(i)
\]
表示语义匹配；

\[
R_T(i)
\]
表示时间匹配；

\[
R_G(i)
\]
表示目标相关；

\[
R_\Psi(i)
\]
表示自我相关。

更加一般地可以写成对数形式：
\[
\log R_i
=
\alpha_S s_S
+
\alpha_W s_W
+
\alpha_T s_T
+
\alpha_G s_G
+
\alpha_\Psi s_\Psi.
\]

这样 Episode Recall 就成为
\[
\boxed{
\mathrm{RelationalEpisodeResonance},
}
\]
而不是单一向量余弦相似度。

不同查询可以动态改变这些权重。例如寻找物体时提高 What 匹配，重新进入过去场景时提高 Where 匹配，反思个人经历时提高 Self 匹配。

\section{E到h：情景召回如何重新进入当前认知}

召回得到的 Episode 不应该直接覆盖当前状态。

设检索结果为
\[
\hat E
=
\sum_i \alpha_iE_i.
\]

首先通过重建器
\[
\mathcal R_E
\]
恢复需要重新激活的状态：
\[
\hat X_W,
\hat X_S,
\hat A,
\hat R
=
\mathcal R_E(\hat E).
\]

然后作为增量注入当前状态：
\[
X_W'
=
X_W
+
g_W\hat X_W,
\]
\[
X_S'
=
X_S
+
g_S\hat X_S,
\]
其中 $g_W,g_S$ 是召回增益。

因此
\[
E\rightarrow h
\]
并不是
\[
\mathrm{Memory}
\rightarrow
\mathrm{ReplaceCurrentState},
\]
而是
\[
\boxed{
\mathrm{Memory}
\rightarrow
\mathrm{BiasCurrentDynamics}.
}
\]

过去经历改变当前状态向哪里演化，但现实输入仍然拥有最终约束权。

\section{从Episode到结构记忆Theta}

长期学习的关键不是无限保留所有 Episode，而是从大量 Episode 中发现稳定关系。

设某类经历集合为
\[
\mathcal E_C
=
\{
E_1,E_2,\ldots,E_n
\}.
\]

如果多个 Episode 反复支持相同结构
\[
C,
\]
则可以定义结构证据：
\[
S_C
=
\sum_i
m_i
\,\operatorname{support}(E_i,C).
\]

当
\[
S_C>\theta_\Theta
\]
并且经过足够现实验证以后，该关系可以进入慢结构：
\[
\boxed{
E\rightarrow\Theta.
}
\]

在 GPU 软件系统中，$\Theta$ 可以由以下对象共同组成：
\[
\Theta
=
\{
\theta_W,
\theta_S,
\theta_C,
\theta_R,
\theta_D,
\ldots
\},
\]
分别表示 What 动力参数、Where 动力参数、跨通路耦合参数、情景召回规则和其他长期动力结构。

因此软件三相记忆中的 $\Theta$ 不只是语言模型权重，而是
\[
\boxed{
\mathrm{ParametersDefiningFutureDynamics}.
}
\]

\section{E到Theta不必等于全模型梯度更新}

一个重要设计原则是：
\[
\boxed{
\mathrm{Consolidation}
\neq
\mathrm{FullModelTraining}.
}
\]

Episode 的长期巩固可以具有多个层级。

最轻量级的是修改 Episode 索引结构：
\[
\Delta\Theta_{\mathrm{index}}.
\]

其次可以形成新的稳定 Prototype：
\[
\Delta\Theta_{\mathrm{prototype}}.
\]

再进一步可以更新局部 Adapter 或专家参数：
\[
\Delta\Theta_{\mathrm{local}}.
\]

只有长期积累和严格验证以后，才允许改变更深层通用动力结构：
\[
\Delta\Theta_{\mathrm{global}}.
\]

因此应形成：
\[
\boxed{
\mathrm{Index}
\rightarrow
\mathrm{Prototype}
\rightarrow
\mathrm{LocalStructure}
\rightarrow
\mathrm{GlobalStructure}.
}
\]

这是一种逐级减慢的结构学习路径。

\section{Theta到E：新的结构重新解释过去}

当长期结构改变以后，旧 Episode 的意义也可能发生变化。

因此 Soft-HGU 不应把每一个事件永久保存成完整自然语言故事。

更加适合的形式是保存
\[
\boxed{
\mathrm{RelationalTrace}
+
\mathrm{RawEvidenceIndex}
+
\mathrm{Metadata}.
}
\]

召回时利用当前 $\Theta$ 生成解释：
\[
E_i^{interpret}
=
\mathcal I(E_i,\Theta_t).
\]

于是：
\[
\boxed{
\Theta_t
\rightarrow
E_i^{interpret}.
}
\]

同一过去经历在不同成长阶段可能得到不同理解，而原始关系证据仍然保持一定独立性。

这能够避免系统把一次早期错误解释永久固化成不可修改的历史。

\section{Theta到h：长期结构塑造当前动力}

软件实现中的长期参数直接定义当前状态的演化函数：
\[
\dot h
=
F_\Theta(h,u).
\]

因此
\[
\Theta\rightarrow h
\]
并不需要执行显式 memory read。

只要
\[
F_\Theta
\]
已经由过去经历改变，那么任何新的认知状态都会自然受到过去学习的影响。

这可以写为：
\[
\boxed{
\mathrm{StructuralMemory}
=
\mathrm{ChangedFutureDynamics}.
}
\]

因此软件 $\Theta$ 与传统数据库的本质区别是：数据库只有被读取时才影响系统，而 $\Theta$ 始终存在于状态转移规律内部。

\section{从h直接写入Theta必须受到保护}

如果系统允许当前一次状态
\[
h(t)
\]
立即修改大量长期参数，就会把偶然噪声、异常输入甚至攻击性数据快速固化成长期结构。

因此默认关系应为：
\[
\boxed{
h
\rightarrow
E
\rightarrow
\Theta,
}
\]
而不是频繁：
\[
h\rightarrow\Theta.
\]

允许直接写入 $\Theta$ 的快速通道必须满足更严格条件：
\[
Confidence>\theta_C,
\]
\[
Risk<\theta_R,
\]
\[
Validation>\theta_V.
\]

这形成软件版本的
\[
\boxed{
\mathrm{SlowVariableProtection}.
}
\]

长期结构越慢、越靠近身份和通用世界模型，其修改门槛越高。

\section{软件几何记忆中的认知温度}

为了让软件系统能够在多个候选状态之间探索，可以定义有效认知温度
\[
T(t).
\]

例如状态转移可以写成：
\[
dh
=
F_\Theta(h)\,dt
+
\sqrt{2T(t)}
\Sigma\,dW_t.
\]

低温状态意味着：
\[
T\downarrow
\Rightarrow
\text{稳定、确定、容易停留在已有吸引子}.
\]

高温状态意味着：
\[
T\uparrow
\Rightarrow
\text{探索范围扩大、更多候选状态参与竞争}.
\]

但必须强调：
\[
\boxed{
T
\text{ controls exploration, not memory writing}.
}
\]

也就是说：
\[
T\uparrow
\]
不能自动导致
\[
\Theta\text{ 更新}.
\]

正确关系是：
\[
\boxed{
\mathrm{Explore}
\rightarrow
\mathrm{Evaluate}
\rightarrow
\mathrm{Episode}
\rightarrow
\mathrm{Validate}
\rightarrow
\mathrm{Consolidate}.
}
\]

这一分离对于保持长期结构稳定至关重要。

\section{元认知控制器的最小定义}

为了控制认知温度、注意增益、抑制程度和搜索深度，可以定义一个较慢的全局控制器
\[
\boxed{
\mathcal C_{\mathrm{meta}}.
}
\]

它接收：
\[
Z(t),
\qquad
U(t),
\qquad
R(t),
\qquad
G(t),
\]
分别表示全局当前状态、不确定度、残差和目标状态。

输出：
\[
u_{\mathrm{meta}}
=
[
T,
G_{\mathrm{attn}},
I_{\mathrm{inh}},
D_{\mathrm{depth}}
].
\]

其中：

$T$ 控制探索温度；

$G_{\mathrm{attn}}$ 控制相关区域增益；

$I_{\mathrm{inh}}$ 控制无关结构抑制；

$D_{\mathrm{depth}}$ 控制继续认知演化的深度。

在 AgentOS 的完整版本中，这一功能后来由更加完整的自我感知和认知控制闭环承担。但在本章软件实验阶段，只需要先实现这一最小控制器即可。

因此本章不依赖任何后续模块才能成立。

\section{从最小元认知控制器到SPC--TCE}

为了与完整 AgentOS 连接，可以进一步定义两个功能角色。

第一类角色负责观察当前认知系统自身的运行状态：
\[
\boxed{
\mathrm{SelfStateObserver}.
}
\]

它把大量分布式状态压缩为
\[
S_{\mathrm{self}}.
\]

第二类角色根据该状态调节认知动力：
\[
\boxed{
\mathrm{CognitiveDynamicsController}.
}
\]

它输出
\[
T,
G,
I,
D,
\ldots
\]

在完整 AgentOS 中，前者称为
\[
\mathrm{SPC},
\]
即 Self-Perception Compressor；

后者称为
\[
\mathrm{TCE},
\]
即 Temperature and Control Engine。

因此本章中的 $Z(t)$ 与元认知控制器可以被理解为 SPC--TCE 的最小软件原型，而不要求读者预先掌握这两个模块。

\section{情感与内稳态状态如何进入软件Episode}

事件对主体的重要性不能只由语义新颖度决定。

系统可以定义一个功能影响状态
\[
A(t)
=
[
a_1,\ldots,a_k
],
\]
表示当前事件对主体资源、安全、任务紧迫性和其他内部变量产生的影响。

这里 $A(t)$ 只是工程状态，不表示主观情绪。

在形成 Episode 时，只保存与当前事件相关的摘要：
\[
A_i^{episode}
=
P_A(A[t_i^{start}:t_i^{end}]).
\]

于是：
\[
E_i
=
(
K_i,
V_i,
\Delta_i,
A_i^{episode},
M_i
).
\]

未来召回该经历时，
\[
A_i^{episode}
\]
可以重新调制当前认知控制，但不会直接覆盖现实身体状态。

因此：
\[
\boxed{
\mathrm{RememberedImpact}
\neq
\mathrm{CurrentImpact}.
}
\]

这避免过去经历制造虚假的当前现实状态。

\section{Self作为Episode选择而不是Episode存储}

长期持续 Agent 还需要一个比普通目标更加稳定的自我结构。

定义慢自我状态：
\[
\Psi(t),
\]
它表示主体长期稳定的身份、价值和经历解释结构。

在软件情景系统中，$\Psi$ 不应该被理解为另一种 Episode Store。

它主要通过以下方式作用于情景系统：

\[
\Psi
\rightarrow
P(\mathrm{write}\ E_i),
\]
\[
\Psi
\rightarrow
P(\mathrm{recall}\ E_i),
\]
\[
\Psi
\rightarrow
P(E_i\rightarrow\Theta).
\]

也就是说，Self 主要决定：
\[
\boxed{
\text{什么经历被认为是``我的经历''，以及这些经历如何影响未来的我。}
}
\]

因此 $\Psi$ 是跨记忆相态的慢约束，而不是第四种普通记忆。

\section{软件HGU不是向量数据库}

Soft-HGU 的底层存储可以使用 GPU Tensor、CPU 内存、SSD、ANN 索引甚至传统数据库。

但是这些都只是 Storage Backend。

Soft-HGU 本身至少必须包含四种过程：
\[
\boxed{
\mathrm{Bind},
}
\]
\[
\boxed{
\mathrm{Store},
}
\]
\[
\boxed{
\mathrm{Complete},
}
\]
\[
\boxed{
\mathrm{Consolidate}.
}
\]

Bind 将当前分布式状态形成 Episode；

Store 管理 Episode 生命周期；

Complete 根据部分线索恢复完整关系；

Consolidate 将稳定跨情景结构逐步写入 $\Theta$。

因此：
\[
\boxed{
\mathrm{SoftHGU}
\neq
\mathrm{VectorDB}.
}
\]

更加准确地说：
\[
\boxed{
\mathrm{SoftHGU}
=
\mathrm{EpisodicDynamicsLayer}
+
\mathrm{StorageBackend}.
}
\]

数据库只是其中最下层的一种实现选择。

\section{GPU分层存储与记忆相态不能混淆}

实际工程中可以让 Episode 按访问速度分层存放：

\[
E_{\mathrm{hot}}
\]
位于 GPU；

\[
E_{\mathrm{warm}}
\]
位于 CPU RAM；

\[
E_{\mathrm{cold}}
\]
位于 SSD 或远程存储。

但：
\[
\boxed{
\mathrm{GPU/CPU/SSD}
}
\]
是物理存储层级；

而：
\[
\boxed{
h/E/\Theta
}
\]
是认知动力学相态。

两者不能等同。

例如一个长期不访问的 Episode 即使存在 SSD 上，仍属于 $E$；

而某个当前被召回并重新进入认知状态的 Episode，其激活部分已经重新成为 $h$。

因此：
\[
\boxed{
\mathrm{MemoryPhase}
\neq
\mathrm{StorageTier}.
}
\]

这一分离使软件三相记忆能够自由使用现有硬件缓存层次，而不破坏理论定义。

\section{软件三相记忆的统一方程}

至此，可以将 GPU 上的软件三相记忆写成如下统一系统。

快速工作状态：
\[
dh
=
F
(
h,
\Theta,
E_{\mathrm{active}},
u,
T
)\,dt
+
\Sigma dW_t.
\]

情景形成：
\[
E_i
=
\mathcal B
\left(
h[t_i^{start}:t_i^{end}]
\right),
\]
满足
\[
p_i^{write}>\theta_E.
\]

情景强度：
\[
\tau_E\dot m_i
=
-\lambda_i m_i
+
\alpha_rR_i
+
\alpha_vV_i.
\]

长期巩固：
\[
\tau_\Theta\dot\Theta
=
P_\Theta
(
\{E_i\},
\varepsilon,
V,
\Theta
).
\]

时间尺度满足：
\[
\boxed{
\tau_h
\ll
\tau_E
\ll
\tau_\Theta.
}
\]

由此形成：
\[
\boxed{
h
\leftrightarrow
E
\leftrightarrow
\Theta.
}
\]

\section{软件实现中的三种Agent-CSO承载形式}

从认知结构对象的角度看，同一个 Agent-CSO 可以在三个相态中具有不同实现。

当前活动态：
\[
C
\mapsto
C_h
=
\mathrm{DistributedActivationPattern}.
\]

情景态：
\[
C
\mapsto
C_E
=
\mathrm{RelationalEpisode}.
\]

结构态：
\[
C
\mapsto
C_\Theta
=
\mathrm{GenerativeParameterStructure}.
\]

因此：
\[
\boxed{
C_h
\neq
C_E
\neq
C_\Theta,
}
\]
但三者可以保持同一认知结构的语义连续性。

由此得到：
\[
\boxed{
\mathrm{MemoryConsolidation}
=
\mathrm{CarrierTransformationOfAgentCSO}.
}
\]

软件三相记忆不是复制同一个向量三次，而是让同一认知结构以越来越慢、越来越稳定的方式存在。

\section{三相记忆的软件频谱解释}

如果把 Agent-CSO 按时间尺度表示为
\[
\rho_{\mathrm{CSO}}(\omega,t),
\]
则 $h$ 主要对应高频活动区：
\[
\rho_h(\omega).
\]

$E$ 对应中频经历区：
\[
\rho_E(\omega).
\]

$\Theta$ 对应低频稳定结构区：
\[
\rho_\Theta(\omega).
\]

于是：
\[
h\rightarrow E
\]
表示结构从快速活动迁向较慢经历；

而
\[
E\rightarrow\Theta
\]
表示经历继续向慢生成结构迁移。

因此：
\[
\boxed{
\mathrm{MemoryConsolidation}
=
\mathrm{CSOSpectralMigration}.
}
\]

召回则执行相反方向：
\[
\Theta/E
\rightarrow
h.
\]

这使软件三相记忆和此前的多尺度智能理论在同一数学语言中统一起来。

\section{Soft-HGU与未来物理HGU的同构关系}

软件 Soft-HGU 与未来可能的物理情景绑定系统不应该被设计成两套完全不同的认知架构。

设软件状态空间为
\[
\mathcal X_S,
\]
物理状态空间为
\[
\mathcal X_H.
\]

若存在映射
\[
\Phi_H:
\mathcal X_S
\rightarrow
\mathcal X_H,
\]
则理想的数字孪生关系要求：
\[
\boxed{
\Phi_H
\left(
D_S^t(x)
\right)
\approx
D_H^t
\left(
\Phi_H(x)
\right).
}
\]

即先在软件中演化再映射到硬件，与先映射到硬件再由物理动力演化，其宏观结果应该近似一致。

对情景绑定同样要求：
\[
\Phi_E
\left(
B_S(W,S,T)
\right)
\approx
B_H
\left(
\Phi_W(W),
\Phi_S(S),
T
\right).
\]

因此软件系统并不是 ``临时替代品''，而可以成为未来物理系统的动力学数字孪生、训练环境和验证平台。

\section{软件实现比直接制造物理芯片更适合作为第一验证阶段}

直接制造可塑物理情景记忆芯片同时涉及器件、封装、连接、噪声、制造良率和长期稳定性等大量变量。

如果认知动力学本身尚未经过实验验证，那么过早进入硬件阶段将很难区分失败究竟来自：

\[
\mathrm{TheoryError}
\]
还是
\[
\mathrm{HardwareError}.
\]

因此更加合理的研究顺序为：
\[
\boxed{
\mathrm{SoftDynamics}
\rightarrow
\mathrm{CognitiveValidation}
\rightarrow
\mathrm{HardwarePhysicalization}.
}
\]

第一阶段先在 GPU 上验证 What--Where 分离是否提高组合泛化能力；

第二阶段验证 Soft-HGU 是否能够实现一次性事件绑定和模式补全；

第三阶段验证
\[
h\leftrightarrow E\leftrightarrow\Theta
\]
是否能够形成长期记忆连续性；

第四阶段再加入 Self、Goal、Affect 和认知控制；

第五阶段才考虑将稳定的动力学机制物理化。

\section{软件实现的第一阶段：双通路验证}

第一阶段只需要构造：
\[
X_W(t),
\qquad
X_S(t),
\]
以及有限的跨通路耦合。

实验目标不是立即建立完整长期 Agent，而是验证：

同一 What 能否在不同 Where 中保持稳定身份；

同一 Where 能否容纳不同 What；

对象换位置以后是否减少重新学习；

空间变化是否不会破坏对象语义；

What 和 Where 是否可以被独立查询和重组。

如果这些能力不能成立，则没有必要继续构造复杂 HGU。

\section{软件实现的第二阶段：Soft-HGU验证}

第二阶段加入：
\[
E_i=(K_i,V_i,M_i).
\]

主要测试：

Where $\rightarrow$ What；

What $\rightarrow$ Where；

Where--When $\rightarrow$ Event；

部分线索 $\rightarrow$ 完整 Episode；

一次经历后的即时召回；

多个相似 Episode 的区分；

情景随时间衰减；

重要经历的稳定保持。

此时尚不需要长期改变 $\Theta$。

\section{软件实现的第三阶段：完整三相循环}

第三阶段正式加入：
\[
E\rightarrow\Theta
\]
以及
\[
\Theta\rightarrow E/h.
\]

主要测试：

重复经历是否形成结构；

长期结构是否减少未来推理成本；

旧 Episode 是否能够被新知识重新解释；

错误 Episode 是否能够衰减而不污染 $\Theta$；

结构学习是否比 Episode 写入明显更慢。

如果上述机制成立，则系统已经具有一个最小的软件三相记忆生命周期。

\section{软件实现的第四阶段：接入完整AgentOS认知控制}

三相记忆自身成立以后，才需要加入更完整的 AgentOS 控制变量。

完整闭环可以写成：
\[
\mathrm{World}
\rightarrow
\mathrm{BodyRuntime}
\rightarrow
\mathrm{CurrentRepresentation}
\rightarrow
h.
\]

随后自状态观察器产生：
\[
S_{\mathrm{self}}.
\]

内稳态与功能影响整合形成：
\[
A_{\mathrm{homeo}}.
\]

认知控制器根据这些状态调节：
\[
T,
\quad
Attention,
\quad
Inhibition,
\quad
Depth.
\]

然后：
\[
h
\rightarrow
E
\rightarrow
\Theta
\rightarrow
\Psi
\rightarrow
\Omega.
\]

慢结构再反向影响未来：
\[
Goal,
\quad
EpisodeWrite,
\quad
Recall,
\quad
Attention,
\quad
Action.
\]

最终形成：
\[
\boxed{
\mathrm{Reality}
\rightarrow
h
\rightarrow
E
\rightarrow
\Theta
\rightarrow
\Psi
\rightarrow
\Omega
\rightarrow
\mathrm{FutureAction}
\rightarrow
\mathrm{Reality}'.
}
\]

\section{软件与物理三相记忆的统一定义}

经过上述构造可以看到，软件实现和物理实现并不是两套三相记忆理论。

软件实现中：
\[
h_S
=
\mathrm{TensorDynamicState},
\]
\[
E_S
=
\mathrm{RelationalEpisodeStructure},
\]
\[
\Theta_S
=
\mathrm{ParametersOfDynamics}.
\]

物理实现中：
\[
h_H
=
\mathrm{PhysicalDynamicState},
\]
\[
E_H
=
\mathrm{PhysicalBindingOrMetastableTrace},
\]
\[
\Theta_H
=
\mathrm{PhysicalGeometry}.
\]

两者都满足：
\[
\boxed{
\mathrm{State}
\leftrightarrow
\mathrm{Trajectory}
\leftrightarrow
\mathrm{Generator}.
}
\]

因此它们的根本区别只是：
\[
\boxed{
\mathrm{DifferentCarrier},
}
\]
而不是
\[
\boxed{
\mathrm{DifferentMemoryTheory}.
}
\]

这再次说明，三相记忆真正描述的是认知结构的时间尺度和动力学角色，而不是某种具体存储技术。

\section{本章总结：在GPU上先建立一个可发育的记忆动力系统}

本章给出了三相记忆从物理构想返回现有 GPU 软件环境的一条完整路径。

首先，把当前认知状态表示为分布式动力变量：
\[
h(t).
\]

其次，把当前表示分化为 What 与 Where 两个可以相互耦合但保持结构独立的空间：
\[
X_W(t),
\qquad
X_S(t).
\]

然后，通过 Soft-HGU 将
\[
\mathrm{What}
+
\mathrm{Where}
+
\mathrm{When}
+
\mathrm{Action}
+
\mathrm{Relation}
\]
绑定成 Episode：
\[
E_i.
\]

Episode 经过衰减、召回、强化和验证以后，稳定统计结构逐渐进入：
\[
\Theta.
\]

于是形成：
\[
\boxed{
h
\leftrightarrow
E
\leftrightarrow
\Theta.
}
\]

这一实现并不要求新的物理芯片。GPU 仍然承担全部底层代数运算，但系统的软件本体已经可以被组织成一个具有几何表示、状态流、情景绑定和多时间尺度结构沉降的认知动力系统。

因此可以把这一研究路线概括为：
\[
\boxed{
\mathrm{GPU}
\rightarrow
\mathrm{SoftwareGeometricDynamics}
\rightarrow
\mathrm{What/Where}
\rightarrow
\mathrm{SoftHGU}
\rightarrow
\mathrm{ThreePhaseMemory}
\rightarrow
\mathrm{AgentLifecycle}.
}
\]

更进一步：
\[
\boxed{
\mathrm{ValidatedSoftDynamics}
\rightarrow
\mathrm{PhysicalGeometricRealization}.
}
\]

其最重要的意义并不是证明未来一定需要某一种几何芯片，而是使三相记忆理论第一次获得一个可以直接在现有计算设备上实验的工程载体。

最终，软件三相记忆可以被压缩成三个定义：
\[
\boxed{
h:
\text{当前仍在发生的分布式计算状态},
}
\]
\[
\boxed{
E:
\text{从当前状态中被切分、绑定并跨时间保存的关系轨迹},
}
\]
\[
\boxed{
\Theta:
\text{由大量经历沉降形成、持续改变未来状态演化方式的慢动力结构}.
}
\]

由此，三相记忆的软件实现并不是在传统 LLM 外部附加一个 ``长期记忆数据库''，而是在计算系统内部重新建立：
\[
\boxed{
\text{状态如何成为经历，经历如何成为结构，结构又如何重新塑造未来状态}
}
\]
这一完整动力闭环。